% =============================================================================
%  Section 4 — Analysis per regime
% =============================================================================
\section{Analysis per regime}
\label{sec:regimes}

\subsection{Pistol}
\label{subsec:pistol}

\begin{figure}[H]
\centering
\includegraphics[width=\linewidth]{pistol_utility_winrate.png}
\caption{Pistol win rate conditional on the total
\texttt{utility\_value} purchased by the team, per quintile.}
\label{fig:pistol_utility}
\end{figure}

In pistol rounds, budgets are nearly equal and tightly bounded
(\$800/player), so very few features discriminate: all correlations
stay below $|r| < 0.06$. The only individual lever worth exploiting
is the \textbf{armour $\leftrightarrow$ utility} trade-off on the T
side, where skipping kevlar (\$650) to fund utility raises the win
rate by about 10 points (Figure~\ref{fig:pistol_utility}). The
\textbf{map} contributes a signal of comparable amplitude
(${\sim}8$~points between \texttt{de\_dust2} and \texttt{de\_anubis}).
HLTV statistics remain uninformative: the historical form of a team
does not transfer to pistol rounds, which behave almost like a coin
flip.

\subsection{Post-pistol}
\label{subsec:postpistol}

\begin{figure}[H]
\centering
\includegraphics[width=0.45\linewidth]{post_pistol_money_scatter.png}
\caption{Post-pistol: CT \texttt{money\_total} vs T
\texttt{money\_total}, coloured by pistol outcome. The asymmetry is
sharp — pistol winner $\approx$ full buy, pistol loser $\approx$
eco/force.}
\label{fig:postpistol_money}
\end{figure}

The pistol outcome creates a structural \textbf{economic asymmetry}
between the winning team, which full-buys, and the losing team, whose
limited cash only allows an eco or a force-buy
(Figure~\ref{fig:postpistol_money}). This asymmetry translates into
\textbf{contrasted buy decisions} on the losing side: a near-systematic
force-buy on the CT side after a lost pistol (78~\%), versus only
48~\% on the T side. The T decision is actually conditioned by the
\emph{plant} of the bomb in the previous round. Its effect is directly
measurable on the force-buy profitability: without a plant the T
force-buy wins 26~\% of the time; with a plant the $+\$600$ per
player bonus is enough to lift it to 41~\%. This chain
\emph{asymmetry $\rightarrow$ decision $\rightarrow$ outcome} gives
post-pistol its exploitable signal, strongly correlated with the
pistol outcome.

\subsection{Normal}
\label{subsec:normal}

\begin{figure}[H]
\centering
\includegraphics[width=0.6\linewidth]{normal_delta_eco_winrate.png}
\caption{CT win rate on normal rounds conditional on $\Delta$
\texttt{equipment\_value} (CT $-$ T), per quintile.}
\label{fig:normal_delta_equip}
\end{figure}

On normal rounds, all buy types coexist (eco, half-buy, force-buy,
full-buy) and the dominant signal is the economic gap measured
directly by $\Delta$ \texttt{equipment\_value}: its lowest quintile
corresponds to a win rate around 24~\% and its top quintile to 77~\%
(Figure~\ref{fig:normal_delta_equip}). At comparable economy (mutual
full-buy with \texttt{equipment\_value} $>$ 20\,k on both sides), two
secondary signals remain. The $\Delta$ \texttt{awp\_count} adds about
5 points of win rate per extra AWP; the map-specific HLTV rating
still separates teams but with a moderate amplitude (win rate
bounded in $[47\,\%, 55\,\%]$ between extreme quintiles).

% =============================================================================
